% -*- coding: utf-8 -*-
% !TEX program = xelatex
\documentclass[plain,noanswer,medmath]{../../template/jnuexam}
\usepackage{../../template/kysx}

\begin{document}


\examtitle{1990}{1}


\makepart{填空题}{本题满分15分，每小题3分}


\begin{problem}
过点$M(1,2,-1)$且与直线
$\begin{cases}
  x=-t+2, \\
  y=3t-4, \\
  z=t-1
\end{cases}$垂直的平面方程是 \fillin{}.
\end{problem}


\begin{problem}
设$a$为非零常数，则$\lim_{x\to\infty}\,\left(\frac{x+a}{x-a}\right)^x=$ \fillin{}.
\end{problem}


\begin{problem}
设函数$f(x)=\begin{cases}
 1, &  |x|\le 1, \\
 0, &  |x|>1, \\
\end{cases}$则$f[f(x)]=$ \fillin{}.
\end{problem}


\begin{problem}
积分$\int_0^2\dx\int_x^2\e^{-y^2}\dy$的值等于 \fillin{}.
\end{problem}


\begin{problem}
已知向量组$\alpha_1=(1,2,3,4)$，$\alpha_2=(2,3,4,5)$，$\alpha_3=(3,4,5,6)$，$\alpha_4=(4,5,6,7)$，
则该向量组的秩是\fillin{}．　　
\end{problem}


\makepart{选择题}{本题满分15分，每小题3分}


\begin{problem}
设$f(x)$是连续函数，且$F(x)=\int_x^{\e^{-x}}\!\!f(t)\dt$，则$F'(x)$等于\pickout{}
\begin{abcd}
\item $-\e^{-x} f(\e^{-x}) - f(x)$．
\item $-\e^{-x} f(\e^{-x}) + f(x)$．
\item $\e^{-x} f(\e^{-x}) - f(x)$．
\item $\e^{-x} f(\e^{-x}) + f(x)$．
\end{abcd}
\end{problem}


\begin{problem}
已知函数$f(x)$具有任意阶导数，且$f'(x)=[f(x)]^2$，
则当$n$为大于$2$的正整数时，$f(x)$的$n$阶导数$f^{(n)}(x)$是\pickout{}
\begin{abcd}
\item $n![f(x)]^{n+1}$．
\item $n[f(x)]^{n+1}$．
\item $[f(x)]^{2n}$．
\item $n![f(x)]^{2n}$．
\end{abcd}
\end{problem}


\begin{problem}
设$\alpha$为常数，则级数$\sum_{n=1}^{\infty}
\left[\frac{\sin n\alpha}{n^2}-\frac{1}{\sqrt{n}}\right]$\pickout{}
\begin{abcd}
\item 绝对收敛．
\item 条件收敛．
\item 发散．
\item 收敛性与$\alpha$的取值有关．
\end{abcd}
\end{problem}


\begin{problem}
已知$f(x)$在$x=0$的某个邻域内连续，且$f(0)=0$，$\lim_{x\to 0}\frac{f(x)}{1-\cos x}=2$，
则在点$x=0$处$f(x)$\pickout{}
\begin{abcd}
\item 不可导．
\item 可导，且$f'(0)\ne 0$．
\item 取得极大值．
\item 取得极小值．
\end{abcd}
\end{problem}


\begin{problem}
已知$\beta_1$, $\beta_2$是非齐次线性方程组$AX=b$的两个不同的解，
$\alpha_1$, $\alpha_2$是对应齐次线性方程组$AX=0$的基础解系，
$k_1$, $k_2$为任意常数，则方程组$AX=b$的通解（一般解）必是\pickout{}
\begin{abcd}
\item $k_1\alpha_1+ k_2(\alpha_1+\alpha_2) + \frac{\beta_1-\beta_2}{2}$．
\item $k_1\alpha_1+ k_2(\alpha_1-\alpha_2) + \frac{\beta_1+\beta_2}{2}$．
\item $k_1\alpha_1+ k_2(\beta_1+\beta_2) + \frac{\beta_1-\beta_2}{2}$．
\item $k_1\alpha_1+ k_2(\beta_1-\beta_2) + \frac{\beta_1+\beta_2}{2}$．
\end{abcd}
\end{problem}


\makepart{计算题}{本题满分15分，每小题5分}


\begin{problem}
求$\int_0^1\frac{\ln(1+x)}{(2-x)^2}\dx$．
\end{problem}


\begin{problem}
设$z=f(2x-y,y\sin x)$，其中$f(u,v)$具有连续的二阶偏导数，求$\frac{\pd^2z}{\pdx\pd y}$．
\end{problem}

        
\begin{problem}
求微分方程$y''+4y'+4y=\e^{-2x}$的通解（一般解）．
\end{problem}


\makepart{}{本题满分6分}

\begin{problem}
求幂级数$\sum_{n=0}^{\infty}(2n+1)x^n$的收敛域，并求其和函数．
\end{problem}


\makepart{}{本题满分8分}

\begin{problem}
求曲面积分$I=\iint_S yz\dx\dx+2\dx\dy$，其中$S$是球面$x^2+y^2+z^2=4$外侧在$z\ge 0$的部分．
\end{problem}


\makepart{}{本题满分7分}

\begin{problem}
设不恒为常数的函数$f(x)$在闭区间$[a,b]$上连续，在开区间$(a,b)$内可导，
且$f(a)=f(b)$．证明在$(a,b)$内至少存在一点$\xi$，使得$f'(\xi)>0$．
\end{problem}


\makepart{}{本题满分6分}

\begin{problem}
设四阶矩阵
$B=\begin{pmatrix}
 1 & -1 &  0 &  0 \\
 0 &  1 & -1 &  0 \\
 0 &  0 &  1 & -1 \\
 0 &  0 &  0 &  1 \\
\end{pmatrix}$，
$C=\begin{pmatrix}
 2 & 1 & 3 & 4 \\
 0 & 2 & 1 & 3 \\
 0 & 0 & 2 & 1 \\
 0 & 0 & 0 & 2 \\
\end{pmatrix}$，
且矩阵$A$满足关系式\newline$A(E-C^{-1}B)'C'=E$，其中$E$为四阶单位矩阵，
$C^{-1}$表示$C$的逆矩阵，$C'$表示$C$的转置矩阵，将上述关系式化简并求矩阵$A$．
\end{problem}


\makepart{}{本题满分8分}

\begin{problem}
求一个正交变换化二次型$f=x_1^2+4x_2^2+4x_3^2-4x_1x_2+4x_1x_3-8x_2x_3$成标准形．
\end{problem}


\makepart{}{本题满分8分}

\begin{problem}
质点$P$沿着以$AB$为直径的圆周，从点$A(1,2)$运动到点$B(3,4)$的过程中受变力$\vec{F}$作用(见图)，
$\vec{F}$的大小等于点$P$与原点$O$之间的距离，其方向垂直于线段$OP$且与$y$轴正向的夹角小于$\frac{\pi}{2}$．
求变力$\vec{F}$对质点$P$所作的功．
\par\centering\includegraphics[width=0.5\linewidth]{1990A}
\end{problem}


\makepart{填空题}{本题满分6分，每小题2分}


\begin{problem}
已知随机变量$X$的概率密度函数$f(x)=\frac{1}{2}\e^{-|x|}$，$-\infty<x<+\infty$，
则$X$的概率分布函数$F(x)=$ \fillin{}.
\end{problem}


\begin{problem}
设随机事件$A$，$B$及其和事件$A\cup B$的概率分别是$0.4$, $0.3$和$0.6$．
若$\overline{B}$表示$B$的对立事件，那么积事件$A\overline{B}$的概率$P(A\overline{B})=$ \fillin{}.
\end{problem}


\begin{problem}
已知离散型随机变量$X$服从参数为$2$的泊松（Poisson）分布，即$P\{X=k\}=\frac{2^k\e^{-2}}{k!}$,
$k=0,1,2,\cdots$，则随机变量$Z=3X-2$的数学期望$E(Z)=$ \fillin{}.
\end{problem}


\makepart{}{本题满分6分}

\begin{problem}
设二维随机变量$(X,Y)$在区域$D:\, 0<x<1,\, |y|<x$内服从均匀分布，
求关于$X$的边缘概率密度函数及随机变量$Z=2X+1$的方差$D(Z)$．
\end{problem}


\end{document}
